% Source PDF: source/普通高考/2012/2012广东文.pdf, page 1, question 9.
% pdfimages -list found no embedded image; the program flowchart is vector artwork.
% Node-edge list: 开始 -> 输入 n -> i=1,s=1 -> i<n?;
% 是 -> s=s×i -> i=i+2 -> left return to the connector above i<n?;
% 否 -> 输出 s -> 结束. No update-to-output edge is present.
% All borders/connectors are solid; labels 是/否 are beside their corresponding branches.
% Every node uses local zero paragraph indentation and centered text.

\begin{tikzpicture}[
  >=Stealth,
  x=1cm,
  y=0.88cm,
  line width=0.45pt,
  line cap=round,
  line join=round,
  font=\normalsize,
  flowtext/.style={anchor=center,align=center,inner sep=0pt,
    execute at begin node={\setlength{\parindent}{0pt}}},
  flowbox/.style={flowtext,draw,minimum height=0.68cm,
    inner xsep=4pt,inner ysep=2pt},
  terminal/.style={flowbox,rounded corners=0.17cm,minimum width=1.35cm},
  io/.style={flowbox,shape=trapezium,trapezium left angle=72,
    trapezium right angle=108,minimum height=0.76cm},
  process/.style={flowbox},
  decision/.style={flowbox,diamond,aspect=2.8,minimum width=3.55cm,
    minimum height=1.05cm},
  branchlabel/.style={flowtext,inner sep=0pt},
]
  \path[use as bounding box] (-3.55,0.10) rectangle (4.85,10.15);

  \node[terminal] (start) at (0,9.65) {开始};
  \node[io,minimum width=2.05cm] (input) at (0,8.35) {输入 $n$};
  \node[process,minimum width=3.45cm] (init) at (0,6.95) {$i=1,s=1$};
  \node[decision] (test) at (0,5.45) {$i<n$};
  \node[process,minimum width=2.85cm,minimum height=0.88cm]
    (update) at (0,3.95) {$s=s\times i$};
  \node[process,minimum width=2.65cm] (increment) at (0,2.50)
    {$i=i+2$};
  \node[io,minimum width=2.10cm] (output) at (3.25,3.85) {输出 $s$};
  \node[terminal] (finish) at (3.25,2.40) {结束};

  % Main path and the affirmative update loop.
  \draw[->] (start.south) -- (input.north);
  \draw[->] (input.south) -- (init.north);
  \draw[->] (init.south) -- (test.north);
  \draw[->] (test.south) -- node[branchlabel,right] {是} (update.north);
  \draw[->] (update.south) -- (increment.north);

  % Negative branch is independent from the update loop.
  \draw[->] (test.east) -- node[branchlabel,above] {否}
    (3.25,5.45) -- (3.25,4.50) -- (output.north);
  \draw[->] (output.south) -- (finish.north);

  % Return from i=i+2 enters the main stem above the test, without resetting i.
  \coordinate (loop-bottom) at (0,1.35);
  \coordinate (loop-left) at (-3.05,1.35);
  \coordinate (loop-entry) at (-3.05,6.25);
  \draw[->] (increment.south) -- (loop-bottom) -- (loop-left)
    -- (loop-entry) -- (0,6.25);
\end{tikzpicture}
